\section{Статус результатов}

В рамках работы был выполнен полный экспериментальный цикл: обучение MRVF-адаптеров, генерация кандидатных ответов, автоматическая попарная оценка и построение таблиц лидеров. Были проведены эксперименты для двух базовых моделей: Qwen3-1.7B и Qwen3-4B.

Первые полные прогоны показали работоспособность программного конвейера, но также выявили проблему в шаблоне генерации: часть ответов содержала объяснения, повторяющиеся фрагменты или не была отделена от рассуждения как короткая финальная шутка. Поскольку автоматическая оценка сравнивает именно сгенерированный текст, такая ошибка формата существенно искажает итоговое качество. После этого шаблон генерации был исправлен: модель должна возвращать только финальную шутку без объяснений. Текущие финальные прогоны с исправленным шаблоном ещё находятся в процессе завершения, поэтому численные значения в таблицах ниже оставлены как заполняемые поля.

Таким образом, на данном этапе зафиксированы структура оценки, объём сравнений и способ агрегации результатов. Окончательные значения доли побед и оценок Брэдли--Терри должны быть внесены после завершения последних прогонов.

\section{Диагностика обучения}

Во время обучения отслеживались значения общей функции потерь, слагаемого для траекторий рассуждения, прямого эталонного слагаемого и средней многореференсной логарифмической массы. Эти метрики позволяют проверить, что обучение действительно оптимизирует внутреннюю MRVF-цель, даже если внешняя оценка юмора может вести себя иначе.

Итоговые диагностические значения для двух обучающих запусков приведены в Таблице~\ref{tab:training-diagnostics}. Значения должны быть заполнены по последним завершённым запускам Weights \& Biases.

\begin{table}[h]
    \caption{Диагностические метрики обучения MRVF}
    \label{tab:training-diagnostics}
    \centering
    \begin{tabular}{lrr}
        \hline
        Метрика                                                  & Qwen3-1.7B & Qwen3-4B \\
        \hline
        Число завершённых шагов                                  & TBD        & TBD      \\
        Итоговая общая потеря                                    & TBD        & TBD      \\
        Итоговая потеря траекторий \(\mathcal{L}_{z}\)           & TBD        & TBD      \\
        Итоговая эталонная потеря \(\mathcal{L}_{\mathrm{ref}}\) & TBD        & TBD      \\
        Средняя \(\log\)-масса                                   & TBD        & TBD      \\
        Средняя длина рассуждения                                & TBD        & TBD      \\
        Доля принудительно закрытых рассуждений                  & TBD        & TBD      \\
        \hline
    \end{tabular}
\end{table}

Эти значения не являются прямой метрикой качества юмора. Их роль состоит в диагностике обучения: если внутренняя \(\log\)-масса растёт, а внешняя оценка не улучшается, это указывает на расхождение между суррогатной обучающей целью и оценкой сгенерированного текста.

\section{Автоматическая оценка Qwen3-1.7B}

Для группы Qwen3-1.7B сравнивались четыре варианта модели: базовая модель без рассуждения, базовая модель с рассуждением, MRVF-модель с рассуждением и MRVF-модель без рассуждения. После фильтрации общих непустых ответов в предварительном прогоне использовалось 184 запроса, что даёт \(184\cdot {4\choose 2}=1104\) попарных сравнений.

Итоговая таблица лидеров для этой группы должна быть заполнена после завершения финального прогона с исправленным шаблоном генерации. Структура таблицы приведена в Таблице~\ref{tab:leaderboard-17b}.

\begin{table}[h]
    \caption{Таблица лидеров для группы Qwen3-1.7B}
    \label{tab:leaderboard-17b}
    \centering
    \begin{tabular}{p{7.2cm}rrrr}
        \hline
        Модель                                                 & BT-оценка & Победы & Поражения & Доля побед \\
        \hline
        \texttt{qwen3-17b-base-nothink-r5eval}                 & TBD       & TBD    & TBD       & TBD        \\
        \texttt{qwen3-17b-base-think-r5eval}                   & TBD       & TBD    & TBD       & TBD        \\
        \texttt{qwen3-17b-mrvf-free-think-r5-budgeted}         & TBD       & TBD    & TBD       & TBD        \\
        \texttt{qwen3-17b-mrvf-free-think-r5-budgeted-nothink} & TBD       & TBD    & TBD       & TBD        \\
        \hline
    \end{tabular}
\end{table}

После заполнения Таблицы~\ref{tab:leaderboard-17b} основной анализ должен отвечать на три вопроса. Во-первых, улучшает ли режим рассуждения базовую модель по сравнению с прямой генерацией. Во-вторых, выигрывает ли MRVF-модель у соответствующей базовой модели в том же режиме вывода. В-третьих, сохраняется ли эффект обученного адаптера, если на этапе генерации отключить явное рассуждение.

\section{Автоматическая оценка Qwen3-4B}

Для группы Qwen3-4B использовался тот же дизайн сравнения: четыре варианта модели и все попарные сравнения между ними. В предварительном прогоне оценка проводилась на 100 запросах, что даёт \(100\cdot {4\choose 2}=600\) попарных сравнений.

Итоговая таблица лидеров для этой группы должна быть заполнена после завершения финального прогона с исправленным шаблоном генерации. Структура таблицы приведена в Таблице~\ref{tab:leaderboard-4b}.

\begin{table}[h]
    \caption{Таблица лидеров для группы Qwen3-4B}
    \label{tab:leaderboard-4b}
    \centering
    \begin{tabular}{p{7.2cm}rrrr}
        \hline
        Модель                                     & BT-оценка & Победы & Поражения & Доля побед \\
        \hline
        \texttt{qwen3-4b-base-nothink}             & TBD       & TBD    & TBD       & TBD        \\
        \texttt{qwen3-4b-base-think}               & TBD       & TBD    & TBD       & TBD        \\
        \texttt{qwen3-4b-mrvf-r1-budgeted-think}   & TBD       & TBD    & TBD       & TBD        \\
        \texttt{qwen3-4b-mrvf-r1-budgeted-nothink} & TBD       & TBD    & TBD       & TBD        \\
        \hline
    \end{tabular}
\end{table}

Сравнение Qwen3-1.7B и Qwen3-4B важно не только как сравнение моделей разного размера. Оно также показывает, насколько чувствителен MRVF-подход к качеству базовой модели. Более крупная модель может иметь более устойчивые траектории рассуждения и более качественные исходные шутки, но при этом требует более дорогого обучения и генерации.

\section{Сравнение режимов вывода}

Для интерпретации результатов полезно рассматривать не только общий рейтинг, но и отдельные разности между близкими вариантами. Таблица~\ref{tab:pairwise-summary} предназначена для краткого анализа основных сравнений после заполнения финальных результатов.

\begin{table}[h]
    \caption{Ключевые сравнения между режимами вывода}
    \label{tab:pairwise-summary}
    \centering
    \begin{tabular}{p{5.2cm}p{5.2cm}p{3.2cm}}
        \hline
        Сравнение                          & Вопрос                                          & Итог \\
        \hline
        Base think против Base no-think    & Помогает ли рассуждение без обучения?           & TBD  \\
        MRVF think против Base think       & Помогает ли MRVF при том же режиме вывода?      & TBD  \\
        MRVF no-think против Base no-think & Сохраняется ли эффект адаптера без рассуждения? & TBD  \\
        MRVF think против MRVF no-think    & Нужно ли рассуждение после MRVF-обучения?       & TBD  \\
        Qwen3-4B против Qwen3-1.7B         & Даёт ли больший размер модели преимущество?     & TBD  \\
        \hline
    \end{tabular}
\end{table}

Эта таблица нужна для того, чтобы не сводить анализ к одному ранжированному списку. Например, если модель с рассуждением проигрывает модели без рассуждения, это может означать, что для коротких шуток длинный промежуточный вывод ухудшает краткость и неожиданность ответа. Если MRVF-модель улучшает внутреннюю \(\log\)-массу, но проигрывает во внешней оценке, это указывает на расхождение между суррогатной обучающей целью и критерием модели-оценщика.

\section{Предварительные наблюдения}

Даже до финального заполнения таблиц можно зафиксировать несколько наблюдений, полученных в ходе отладочных прогонов.

Во-первых, MRVF-обучение технически реализуемо для моделей Qwen3-1.7B и Qwen3-4B с использованием LoRA-адаптеров. Обучение не требует отдельной модели награды: сигнал получается из условного правдоподобия эталонных шуток при сгенерированной траектории рассуждения.

Во-вторых, качество итоговой оценки оказалось чувствительным к формату генерации. Ошибка в шаблоне вывода приводила к тому, что модели генерировали пояснения или повторяющиеся фрагменты вместо одной короткой финальной шутки. Это особенно важно для юмора: даже хорошая идея может получить низкую оценку, если ответ слишком длинный, объяснительный или не похож на законченную шутку.

В-третьих, предварительные запуски показали, что улучшение внутренней обучающей цели не гарантирует немедленного выигрыша в автоматической попарной оценке. Это соответствует ожидаемому риску MRVF-постановки: логарифмическая масса эталонных ответов является суррогатной целью, а не прямой мерой человеческой смешности.

\section{Интерпретация итоговых результатов}

После завершения финальных прогонов интерпретация должна строиться по следующей схеме.

Если MRVF-модели превосходят соответствующие базовые модели, это будет означать, что многореференсная логарифмическая цель действительно даёт полезный обучающий сигнал для генерации юмора. В таком случае основной вывод будет состоять в том, что даже неполные корпусные множества эталонных шуток могут использоваться как слабый сигнал для улучшения генерации без отдельной модели награды.

Если MRVF-модели уступают базовым моделям, такой результат не следует интерпретировать как чисто техническую неудачу. Он будет означать, что текущая суррогатная цель, короткий бюджет обучения или малый размер группы траекторий недостаточны для улучшения внешнего качества. В этом случае важным результатом работы становится выявление расхождения между оптимизацией многореференсной \(\log\)-массы и автоматической оценкой юмора.

Если модели с рассуждением уступают моделям без рассуждения, это будет указывать на особенность задачи: короткая шутка часто требует компактного финального текста, а длинное рассуждение может ухудшать стиль, краткость или неожиданность ответа. Если же режим рассуждения улучшает качество, это будет подтверждать гипотезу о пользе промежуточного творческого поиска.

\section{Ограничения результатов}

Полученные результаты имеют несколько ограничений.

Первое ограничение связано с оценкой. В работе используется автоматическая модель-оценщик, а не человеческая разметка. Поэтому итоговые доли побед и оценки Брэдли--Терри измеряют качество относительно выбранного протокола автоматической оценки, а не абсолютную человеческую смешность.

Второе ограничение связано с неполнотой эталонных множеств. Набор \(Y^{*}(x)\) содержит только найденные корпусные шутки и не исчерпывает все возможные хорошие ответы на запрос. Следовательно, MRVF-цель может усиливать похожесть на наблюдаемые эталоны, но не обязана напрямую улучшать творческую оригинальность.

Третье ограничение связано с масштабом обучения. Проведённые запуски являются отладочными и используют ограниченное число шагов, малый размер группы траекторий и небольшое число эталонных ответов на запрос. Полноценная проверка метода требует более длинных запусков и абляций по числу траекторий, числу эталонов и коэффициентам функции потерь.

Четвёртое ограничение связано с форматом вывода. Первый полный прогон показал, что результаты существенно зависят от точности шаблона генерации. Поэтому все итоговые выводы должны делаться только по финальным прогонам, в которых модель возвращает одну финальную шутку без объяснения.

\section{Выводы по экспериментальной части}

Экспериментальная часть подтверждает реализуемость полного конвейера: от построения многореференсного набора данных до обучения LoRA-адаптеров, генерации кандидатов, автоматической попарной оценки и построения таблиц лидеров. Главный методический вывод состоит в том, что для юмора важно оценивать не только внутреннюю обучающую цель, но и фактический формат и качество финального ответа.

Окончательные выводы о сравнении MRVF-моделей с базовыми моделями должны быть сделаны после завершения финальных прогонов с исправленным шаблоном генерации и заполнения Таблиц~\ref{tab:leaderboard-17b} и~\ref{tab:leaderboard-4b}.